\documentclass[11pt]{article}
\usepackage[a4paper,margin=1in]{geometry}
\usepackage{amsmath,amssymb,amsthm}
\usepackage{booktabs}
\usepackage{array}
\usepackage{hyperref}

\newtheorem{definition}{Definition}
\newtheorem{theorem}{Theorem}
\newtheorem{nogo}{No-Go}

\newcommand{\dB}{\mathrm{dB}}
\newcommand{\HB}{\mathrm{HB}}
\newcommand{\XiBC}{\Xi_{\mathrm{BC}}}
\newcommand{\XiHecke}{\Xi_{\mathrm{BC,Hecke}}}
\newcommand{\XidB}{\Xi_{\dB}}
\newcommand{\Hzeta}{H(E_\zeta)}
\newcommand{\Estar}{E^\#}

\title{Step 181: de Branges Carrier Pivot}
\author{RH Membrane Adequacy Track}
\date{}

\begin{document}
\maketitle

\section*{Verdict}

\[
\boxed{\mathrm{dB\_verdict}=V_{\dB\_\mathrm{target\_equivalent}}.}
\]

The de Branges carrier can be declared as a typed Hilbert space of entire
functions, but the RH-relevant choice of Hermite-Biehler data is not free
carrier data.  Making it legal, or proving the corresponding de Branges
subspace-chain positivity, is exactly the classical de Branges RH program.
Thus the pivot yields a new typed no-go:
\[
\boxed{\text{de Branges RH-carrier closure target-equivalence / Conrey--Li survival}.}
\]

\section*{Inherited Records Quoted}

\paragraph{Step 145: Burnol/Sonine carrier.}
The inherited Burnol/Sonine carrier is not a de Branges carrier package.  The
record declares
\[
L_a\subset L^2(0,\infty;dt)
\]
as Burnol's extended Sonine space: functions constant on \((0,a)\) whose
cosine transform is again constant on \((0,a)\).  It then defines
\[
H_R=\overline{\operatorname{span}\{
\operatorname{Ran}\Pi_{Y_a}\mathfrak B_{\ell,a}:
\ell\in\mathcal L\setminus\{0\}\}}\subseteq Y_a\subset L_a.
\]
This gives a concrete Burnol/Sonine residual carrier, not an explicit
Hermite-Biehler function \(E\).

\paragraph{Step 153: projected Sonine RKHS and transport.}
The pulled-evaluator record gives, for \(w\notin\{0,1\}\),
\[
[f,Y^a_{w,k}]_{L_a}=M(f)^{(k)}(w), \qquad f\in L_a,
\]
and defines the completed Mellin image \(\mathcal L_a^\Gamma\).  Its
reproducing kernel is recorded by
\[
G(w)=\langle G,K_a^\Gamma(\cdot,w)\rangle_{\mathcal L_a^\Gamma}.
\]
The transport is
\[
T_a:=\mathcal M_\Gamma J_a\mathcal U_\infty^{-1},
\]
and the pulled evaluator is
\[
\mathcal U_\infty(J_a^*Y^a_{w,k})
   =T_a^*\partial_{\overline w}^{\,k}K_a^\Gamma(\cdot,w).
\]
The same record warns that
\[
K_a^\Gamma(\cdot,w)=P_{\mathcal L_a^\Gamma}
K_a^{\Gamma,\mathrm{amb}}(\cdot,w),
\]
so the Sonine projection is load-bearing.  Dropping it gives only an ambient
Hardy-shadow formula.

\paragraph{Step 90 and Step 93: de Branges warning.}
Step 90 states that de Branges spaces supply Hilbert kernels and positivity
candidates, but that the needed phase/kernel positivity is essentially the
hard gate of the de Branges RH program.  Step 93 records the Conrey--Li
survival test: natural de Branges-type positivity conditions associated with
\(\xi(s)\) and \(1/\xi(s)\) fail for zeta and Dirichlet \(L\)-functions, so
any new positivity ansatz must not imply those refuted conditions.

\paragraph{Corpus audit.}
The local files \texttt{anti\_loc/adequacy.tex} and
\texttt{anti\_loc/needles.tex} do not contain an accepted de Branges carrier
package.  They therefore do not supply an inherited \(E_\zeta\), an
\(H(E_\zeta)\)-to-\(L_a\) bridge, or a completed de Branges residual theorem.

\section*{Classical de Branges Carrier}

\begin{definition}[Hermite-Biehler function and de Branges space]
Let \(E\) be an entire function and
\[
E^\#(z):=\overline{E(\overline z)}.
\]
The function \(E\) is Hermite-Biehler when
\[
|E(z)|>|E^\#(z)| \qquad (\operatorname{Im}z>0),
\]
with the usual no-real-zero and growth hypotheses of the de Branges theory.
The de Branges space \(H(E)\) is the Hilbert space of entire functions with
the de Branges norm, and its reproducing kernel is
\[
K_E(z,w)=
\frac{
E(z)\overline{E(w)}-E^\#(z)\overline{E^\#(w)}
}{
2\pi i(\overline w-z)
}.
\]
\end{definition}

The RH-linked entire function is
\[
\xi(s)=\frac12s(s-1)\pi^{-s/2}\Gamma(s/2)\zeta(s),
\qquad
\Xi(t)=\xi\!\left(\frac12+it\right).
\]
A natural de Branges RH carrier would require
\[
E_\zeta(z)=A_\zeta(z)-iB_\zeta(z),\qquad A_\zeta(z)=\Xi(z),
\]
with \(B_\zeta\) a real entire companion such that \(E_\zeta\) is
Hermite-Biehler, or a unitarily equivalent normalization.  This declaration is
conditional: the existence of such \(B_\zeta\), together with the relevant
chain positivity/order statement, is not inherited free data.

\section*{Operators and Parent Residual}

The typed de Branges carrier supplies the following native objects.
\[
\begin{array}{ll}
S_E & \text{canonical symmetric multiplication }F(z)\mapsto zF(z),\\
K_E(\cdot,w) & \text{kernel/evaluator vector in }H(E),\\
P_{\mathrm{chain}}(t) & \text{de Branges subspace-chain projection, if supplied.}
\end{array}
\]
The Burnol commutator
\[
C_\ell=(I-P_\infty)M_{m_\ell}P_\infty
\]
does not canonically transport to \(H(E_\zeta)\) from the inherited records.
A bridge from \(L_a\) or \(\XiBC\) to \(\Hzeta\) would be a separate typed
obligation.

Define the de Branges parent residual \(\XidB\) as the typed defect of the
RH-specific de Branges structure:
\[
\XidB=0
\quad\Longleftrightarrow\quad
\text{\(E_\zeta\) is legal Hermite-Biehler and the required chain condition holds.}
\]
This is a declaration of the residual, not a proof that the residual vanishes.

\section*{Initial Cascade on \(H(E_\zeta)\)}

\begin{center}
\begin{tabular}{p{0.13\linewidth}p{0.25\linewidth}p{0.35\linewidth}p{0.18\linewidth}}
\toprule
Branch & Name & Target & Status\\
\midrule
dB-E & Hermite-Biehler companion legality
& Construct \(E_\zeta=A_\zeta-iB_\zeta\), \(A_\zeta=\Xi\), with \(E_\zeta\) Hermite-Biehler.
& target-equivalent\\
dB-chain & de Branges chain positivity/order
& Prove the chain condition associated with \(E_\zeta\).
& target-equivalent\\
dB-kernel & kernel positivity/source ladder
& Use \(K_E\) matrices as source positivity.
& blocked by Conrey--Li survival and public-shadow no-go\\
dB-bridge & bridge to Burnol/Sonine
& Compare \(H(E_\zeta)\) with the Burnol/Sonine carrier of \(\XiBC\).
& no accepted bridge\\
\bottomrule
\end{tabular}
\end{center}

This cascade is not a new non-CTMT attack route.  It does not reach CTMT as
the primary obstruction because it stops earlier: the RH-specific
Hermite-Biehler and chain data are already target-equivalent.

\section*{No-Go Transfer Audit}

\begin{center}
\begin{tabular}{p{0.34\linewidth}p{0.14\linewidth}p{0.42\linewidth}}
\toprule
No-go & Transfer & Reason\\
\midrule
public-shadow non-promotion & yes
& naive RKHS/de Branges positivity is a public-shadow shortcut unless carrier-complete records are supplied\\
finite-window Calkin blindness & yes
& finite kernel matrices do not decide completed chain positivity\\
auxiliary-GRH smuggling & yes
& real-zero or chain positivity for auxiliary entire functions cannot be assumed\\
incomplete character spectrum support-only & retained
& source-side no-go remains for any attempted source ladder\\
scalar identity not carrier identity & yes
& scalar \(\xi\) identities do not identify \(H(E)\) carriers\\
ledger-relative Hecke non-comparability & retained
& no Hecke-to-de-Branges carrier ledger bridge is supplied\\
Branch C full-carrier ZI-COV(i) foreclosure & retained
& Burnol/Sonine route-local no-go remains, but does not decide \(H(E)\)\\
de Branges RH-carrier closure target-equivalence / Conrey--Li survival & new
& natural de Branges RH positivity/chain closure is target-equivalent or refuted in naive zeta RKHS forms\\
\bottomrule
\end{tabular}
\end{center}

\begin{nogo}[de Branges target-equivalence no-go]
Under the inherited records, closure of the natural RH-linked de Branges
carrier residual \(\XidB\) is target-equivalent to RH and cannot be used as an
independent adequacy route.
\end{nogo}

\begin{proof}
The only RH-relevant de Branges carrier declared in this step is an
\(H(E_\zeta)\) whose real part is \(\Xi\), up to harmless normalization:
\[
E_\zeta=A_\zeta-iB_\zeta,\qquad A_\zeta=\Xi.
\]
For this to be a valid carrier, \(E_\zeta\) must be Hermite-Biehler and must
satisfy the de Branges chain/order condition required by the classical RH
program.  Classical de Branges theory supplies the kernel and ordered
subspace formalism, but it does not make this RH-specific positivity free.
The de Branges RH program identifies exactly these structural conditions as
the hard RH content.

The inherited source audit also blocks the naive kernel-positivity shortcut:
Conrey--Li refutes natural de Branges-type positivity conditions associated
with \(\xi(s)\) and \(1/\xi(s)\).  Therefore an argument that simply promotes
zeta RKHS positivity would violate the inherited public-shadow and
Conrey--Li survival constraints.

Finally, Step 145 and Step 153 give Burnol/Sonine \(L_a\), \(T_a\), and the
projected Sonine kernel \(K_a^\Gamma\), but they do not give an explicit
\(E_\zeta\), an \(H(E_\zeta)\)-to-\(L_a\) bridge, or an equality of residuals.
Thus \(\XidB=0\) is not an independent closure of \(\XiBC\); it is the
classical de Branges RH structural condition restated as a typed residual.
\end{proof}

\section*{Classical Theorems Cited}

\begin{itemize}
\item L. de Branges, \emph{Hilbert Spaces of Entire Functions}, 1968:
Hermite-Biehler spaces, reproducing kernels, and de Branges subspace ordering.
\item de Branges RH program papers, 1986 and later: RH-related formulation by
de Branges chain/positivity conditions for zeta-related entire functions.
\item J. B. Conrey and X.-J. Li, 2000: counterexamples to natural
de Branges-type positivity conditions for zeta and Dirichlet \(L\)-functions.
\end{itemize}

\section*{Nonclaim Boundary}

This step does not prove RH, does not assume RH, and does not close
\(\XiBC\), \(\XiHecke\), or \(\XidB\).  It does not assume the existence of a
Hermite-Biehler \(E_\zeta\) with \(A_\zeta=\Xi\), does not use de Branges
chain positivity as a tool, and does not adopt the naive zeta RKHS positivity
conditions warned against by Conrey--Li.  The result is a typed target-
equivalence no-go for the natural de Branges RH carrier closure.

\end{document}
